\documentclass[aps,prd,a4paper,10pt, twocolumn,amsmath,showpacs,superscriptaddress,nofootinbib,preprintnumbers,notitlepage]{revtex4-2}
\usepackage{dcolumn}
\usepackage{epstopdf}
\usepackage{bm}
\usepackage{braket}
\usepackage{enumitem}
\usepackage{soul}
\usepackage{pifont}
\usepackage[labelfont= bf, skip=5pt, font=small]{caption}
\usepackage{graphicx, wrapfig, subcaption}
\usepackage{array, tabularx, tabulary, booktabs, multirow}
\usepackage{textcomp, cmap, comment, epigraph}
\usepackage[dvipsnames]{xcolor}
\usepackage{orcidlink}
\usepackage[utf8]{inputenc}
\usepackage[spanish, es-tabla]{babel}
\usepackage[version=4]{mhchem}
\usepackage{hyperref}
\usepackage{geometry}
\usepackage{amsfonts, amssymb, amsthm, mathtools} 
\usepackage{amsmath}
\usepackage{float}
\usepackage{placeins}
\usepackage{siunitx}

\geometry{left=15mm, right=15mm, top=16mm, bottom=18mm}
\DeclareMathSizes{12}{16}{12}{10}
\selectlanguage{spanish}

\makeatletter
\newcommand{\fmarki}{*}
\newcommand{\fmarkii}{\ensuremath{\dagger}}
\newcommand{\fmarkiii}{\ensuremath{\ddagger}}  
\def\@fnsymbol#1{{\ifcase#1\or \fmarki\or \fmarkii\or \fmarkiii\else\@ctrerr\fi}}
\makeatother

\renewcommand{\thesection}{\Roman{section}}
\renewcommand{\abstractname}{}

\addto\captionsspanish{\renewcommand{\abstractname}{}}
\def\restrict#1{\raise-.5ex\hbox{\ensuremath|}_{#1}}
\def\sen{\mathop{\mbox{\normalfont sen}}\nolimits}
\graphicspath{ {images/} }
\hypersetup{colorlinks=true,linkcolor=MidnightBlue,citecolor=OliveGreen,urlcolor=RoyalBlue,pdftitle={},pdfauthor={Nicolás Clotta}}

\begin{document}

\title{Cuaderno de Laboratorio 6}

\author{Nicolás Clotta}
\email{nicolas.clotta@gmail.com}
\affiliation{Departamento de Física, Facultad de Ciencias Exactas y Naturales, Universidad de Buenos Aires.}

\date[Fechado: ]{\today}
%\begin{abstract}
%    \noindent
%\end{abstract}

\maketitle

\section{Introducción}
\subsection{Nucleación}
Proceso mediante el cuál se forman defectos en el material debido a la irradiación recibida. Ocurre primero sobre las aglomeraciones de cobre, y tiende formar cúmulos de mayor tamaño a su alrededor. Luego, ocurre sobre las de manganeso. El proceso físico que lleva estos defectos a tomar mayor relevancia frente a los naturales se debe al decaimiento $\beta^+$, en este caso causando que los átomos de hierro se conviertan en átomos de manganeso, aunque también hay una gran prevalencia de conversión atómica por captura electrónica \cite{enwiki:isotopes_of_iron, takahashi2004}.

\subsection{}


\section{Minor loops, paredes de dominio}

La posibilidad de medir la magnetización 

\section{Lake Shore VSM}
\begin{itemize}
\item $m_s$ momento de saturación: A partir de valores del campo entre los cuales el moment magnético es lineal, 

\item $H_{ci}$ coercitividad intrínseca

\item $H_c$ coercitividad

\item $m_r$ remanencia

\item $B_r$ inducción residual

\item $\displaystyle\frac{m_r}{m_s}$ squareness ratio

\item $B\times H_{max}$ max. energy product
\end{itemize}

\section{Dimensiones \& peso muestras}
Calibración según el microscopio e ImageJ: $(0.6658\pm0.0001)~\mathrm{px}/\mathrm{\mu m}$. Se tomó un valor de calibración nuevo por cada imagen por como funciona ImageJ aunque creo que no debería hacer falta simplemente soy un gil.
\subsection{Disco}
\begin{itemize}
\item Peso: $(0.0075\pm0.0001)~\mathrm{g}$.
\item Radio vertical: $(1706.197\pm0.001)~\mathrm{\mu m}$.
\item Radio horizontal: $(1703.206\pm0.001)~\mathrm{\mu m}$.
\item Ancho: $(123.150\pm0.001)~\mathrm{\mu m}$.
\end{itemize}

\subsection{Cuboide}
\begin{itemize}
\item Peso: $(0.0057\pm0.0001)~\mathrm{g}$.
\end{itemize}

\subsection{Barra}
\begin{itemize}
\item Peso: $(0.0615\pm0.0001)~\mathrm{g}$.
\item Largo: $(0.95\pm0.01)~\mathrm{mm}$.
\item Lado A, cara 1: $(985.065\pm0.001)~\mathrm{\mu m}$.
\item Lado B, cara 1 $(900.718\pm0.001)~\mathrm{\mu m}$.
\item Medio, cara 2: $(1095.301\pm0.001)~\mathrm{\mu m}$.

\end{itemize}

\subsection{Prisma}
\begin{itemize}
\item Peso: $(0.0348\pm0.0001)~\mathrm{g}$.
\item Largo: $(0.60\pm0.01)~\mathrm{mm}$.
\item Lado A, cara 1: $(1396.177\pm0.001)~\mathrm{\mu m}$.
\item Medio, cara 1: $(1249.993\pm0.001)~\mathrm{\mu m}$.
\item Lado B, cara 1 $(1193.333\pm0.001)~\mathrm{\mu m}$.
\item Lado A, cara 2: $(908.925\pm0.001)~\mathrm{\mu m}$.
\item Medio, cara 2: $(914.772\pm0.001)~\mathrm{\mu m}$.
\item Lado B, cara 2: $(1006.342\pm0.001)~\mathrm{\mu m}$.
\end{itemize}

\section{Resultados Ajustes}
\subsection{Disco $0^\circ$}

$W^0_F = (29483.712\pm6986.505)$ [ergs/g]

$n_f=(2.209\pm0.208)$




$W^0_r = (6.962\pm0.010)$ [ergs/g]

$n_r=(1.946\pm0.003)$




$H^0_c = (7.810\pm0.011)$ [Oe]

$n_c = (0.946\pm0.003)$

Los exponentes dan casi lo mismo que el paper de Sebastián, creo que están bien.
\subsection{Disco $90^\circ$}
$W^0_F = (97426.811\pm21256.146)$ [ergs/g]

$n_f=(2.374\pm0.063)$



$W^0_r = (135.169\pm2.818)$ [ergs/g]

$n_r=(2.000\pm0.005)$

$H^0_c = (151.619\pm3.161)$ [Oe]

$n_c = (1.000\pm0.005)$

Acá da todo cualquier cosa. Es horrible.

\section{Papers a revisar}

%http://doi.org/10.3390/cryst12081091

\bibliographystyle{plain}
\bibliography{cuaderno6}




\end{document}
